% !TEX program = xelatex
% ============================================================
%  وثيقة انطلاق المشروع — Project Kickoff Document Template
%  نموذج لوثيقة انطلاق المشروع تجمع كل المعلومات الأساسية في مكان واحد
%  Compile with: xelatex main.tex (run twice)
% ============================================================
%  Features:
%    - Project overview, objectives, scope (in/out)
%    - Key deliverables table
%    - Timeline section
%    - Team roles table
%    - Success criteria
%    - Risks placeholder
%    - Communication plan
% ============================================================

\documentclass[11pt,a4paper]{article}

\usepackage[a4paper,margin=2.5cm,top=2.5cm,bottom=2.5cm,headheight=16pt]{geometry}
\usepackage{graphicx}
\usepackage{array}
\usepackage{booktabs}
\usepackage{tabularx}
\usepackage{multirow}
\usepackage{enumitem}
\usepackage{float}
\usepackage{titlesec}
\usepackage{fancyhdr}
\usepackage{setspace}
\usepackage{xcolor}
\usepackage{amsmath}

\usepackage{bidi}
\usepackage[no-math]{fontspec}
\usepackage[quiet]{polyglossia}

\setmainfont[Scale=1]{NotoKufiArabic-Regular.ttf}
\newfontfamily\arabicfont[Script=Arabic,Scale=0.95]{NotoKufiArabic-Regular.ttf}
\newfontfamily\englishfont[Scale=0.95]{TeX Gyre Termes}

\setdefaultlanguage[calendar=gregorian,hijricorrection=1]{arabic}
\setotherlanguage[variant=british]{english}

\usepackage[breaklinks,hidelinks]{hyperref}

\addto\captionsarabic{%
  \renewcommand{\contentsname}{الفهرس}%
  \renewcommand{\figurename}{الشكل}%
  \renewcommand{\tablename}{الجدول}%
}

\titleformat{\section}{\Large\bfseries}{\thesection}{0.7em}{}
\titlespacing*{\section}{0pt}{1.4ex plus 0.4ex}{0.9ex plus 0.3ex}

\pagestyle{fancy}
\fancyhf{}
\fancyhead[R]{\small\itshape وثيقة انطلاق المشروع}
\fancyhead[L]{\small اسم المشروع}
\fancyfoot[C]{\small\thepage}
\renewcommand{\headrulewidth}{0.3pt}

\setlist[itemize]{topsep=0.3em, itemsep=0.15em, parsep=0pt}
\setlist[enumerate]{topsep=0.3em, itemsep=0.15em, parsep=0pt}

\newcolumntype{L}[1]{>{\raggedright\arraybackslash}p{#1}}
\newcolumntype{C}[1]{>{\centering\arraybackslash}p{#1}}

\setstretch{1.3}

\begin{document}

% ============================================================
% TITLE BLOCK
% ============================================================
\begin{center}
{\LARGE\bfseries وثيقة انطلاق المشروع}\\[0.3em]
{\large Project Kickoff Document}\\[1em]
\end{center}

% ============================================================
% PROJECT INFO
% ============================================================
% TODO: املأ بيانات المشروع.
\begin{table}[H]
\centering
\renewcommand{\arraystretch}{1.4}
\begin{tabularx}{\textwidth}{L{3.5cm} X L{3.5cm} X}
\toprule
\textbf{اسم المشروع:} & --- & \textbf{مدير المشروع:} & --- \\
\textbf{اسم الشركة:} & --- & \textbf{الراعي الرئيسي:} & --- \\
\textbf{تاريخ البدء:} & --- & \textbf{تاريخ الانتهاء:} & --- \\
\textbf{الميزانية:} & --- & \textbf{تاريخ الوثيقة:} & --- \\
\bottomrule
\end{tabularx}
\end{table}

\vspace{1em}

% ============================================================
% 1. PROJECT OVERVIEW
% ============================================================
\section{نظرة عامة على المشروع}
\label{sec:overview}

% TODO: اكتب نظرة عامة مختصرة على المشروع.
هذا نص تجريبي للنظرة العامة. يقدم وصفاً مختصراً للمشروع وأهميته والأسباب التي دعت إلى إطلاقه، مع توضيح القيمة المتوقعة للمنظمة والمستفيدين.

% ============================================================
% 2. OBJECTIVES
% ============================================================
\section{أهداف المشروع}
\label{sec:objectives}

% TODO: اكتب أهداف المشروع.
\begin{enumerate}
    \item الهدف الأول: وصف الهدف القابل للقياس.
    \item الهدف الثاني: وصف الهدف القابل للقياس.
    \item الهدف الثالث: وصف الهدف القابل للقياس.
\end{enumerate}

% ============================================================
% 3. SCOPE
% ============================================================
\section{نطاق المشروع}
\label{sec:scope}

\subsection{المشمولات}
% TODO: اذكر ما يشمله المشروع.
\begin{itemize}
    \item البند الأول المشمول في المشروع.
    \item البند الثاني المشمول في المشروع.
    \item البند الثالث المشمول في المشروع.
\end{itemize}

\subsection{غير المشمولات}
% TODO: اذكر ما لا يشمله المشروع.
\begin{itemize}
    \item البند الأول غير المشمول.
    \item البند الثاني غير المشمول.
\end{itemize}

% ============================================================
% 4. KEY DELIVERABLES
% ============================================================
\section{المخرجات الرئيسية}
\label{sec:deliverables}

% TODO: عدّل المخرجات حسب مشروعك.
\begin{table}[H]
\centering
\caption{المخرجات الرئيسية للمشروع}
\renewcommand{\arraystretch}{1.3}
\begin{tabularx}{\textwidth}{C{1cm} L{5cm} L{3.5cm} X}
\toprule
\textbf{رقم} & \textbf{المخرج} & \textbf{تاريخ التسليم} & \textbf{معايير القبول} \\
\midrule
1 & وثيقة المتطلبات & --- & موافقة الراعي الرئيسي \\
2 & التصميم المعماري & --- & مراجعة لجنة التقنية \\
3 & النموذج الأولي & --- & موافقة ممثل المستخدمين \\
4 & النسخة التجريبية & --- & اجتياز اختبارات القبول \\
5 & النسخة النهائية & --- & اجتياز جميع اختبارات الجودة \\
6 & وثائق التشغيل والصيانة & --- & موافقة فريق التشغيل \\
\bottomrule
\end{tabularx}
\end{table}

% ============================================================
% 5. TIMELINE
% ============================================================
\section{الجدول الزمني}
\label{sec:timeline}

% TODO: حدّد المراحل الزمنية الرئيسية.
\begin{table}[H]
\centering
\caption{المراحل الزمنية الرئيسية}
\renewcommand{\arraystretch}{1.3}
\begin{tabularx}{\textwidth}{C{1cm} L{4cm} L{3cm} L{3cm} X}
\toprule
\textbf{رقم} & \textbf{المرحلة} & \textbf{تاريخ البدء} & \textbf{تاريخ الانتهاء} & \textbf{المعالم الرئيسية} \\
\midrule
1 & التحليل والتخطيط & --- & --- & اعتماد المتطلبات \\
2 & التصميم & --- & --- & اعتماد التصميم المعماري \\
3 & التطوير & --- & --- & النموذج الأولي جاهز \\
4 & الاختبار & --- & --- & اجتياز اختبارات القبول \\
5 & الإطلاق والتشغيل & --- & --- & النسخة النهائية منشورة \\
\bottomrule
\end{tabularx}
\end{table}

% ============================================================
% 6. TEAM ROLES
% ============================================================
\section{أدوار فريق المشروع}
\label{sec:team}

% TODO: عدّل أدوار الفريق حسب مشروعك.
\begin{table}[H]
\centering
\caption{أدوار ومسؤوليات فريق المشروع}
\renewcommand{\arraystretch}{1.3}
\begin{tabularx}{\textwidth}{L{3cm} L{3.5cm} X}
\toprule
\textbf{الاسم} & \textbf{الدور} & \textbf{المسؤوليات الرئيسية} \\
\midrule
--- & مدير المشروع & الإشراف العام وإدارة الجدول والميزانية \\
--- & قائد الفريق التقني & قيادة الفريق التقني واتخاذ القرارات المعمارية \\
--- & مطور أول & تطوير المكونات الأساسية ومراجعة الكود \\
--- & مطور & تنفيذ مهام التطوير المحددة \\
--- & مختبر الجودة & إعداد وتنفيذ اختبارات الجودة \\
--- & محلل الأعمال & جمع المتطلبات وتوثيقها \\
--- & مسؤول الاتصال & إدارة التواصل مع أصحاب المصلحة \\
\bottomrule
\end{tabularx}
\end{table}

% ============================================================
% 7. SUCCESS CRITERIA
% ============================================================
\section{معايير النجاح}
\label{sec:success}

% TODO: حدّد معايير النجاح القابلة للقياس.
\begin{enumerate}
    \item تسليم جميع المخرجات في المواعيد المحددة وبالميزانية المعتمدة.
    \item اجتياز النظام لجميع اختبارات الجودة بنسبة لا تقل عن 95\%.
    \item رضا المستخدمين بنسبة لا تقل عن 80\% وفقاً لاستبيان ما بعد الإطلاق.
    \item عدم تجاوز الميزانية المعتمدة بأكثر من 5\%.
\end{enumerate}

% ============================================================
% 8. RISKS
% ============================================================
\section{المخاطر}
\label{sec:risks}

% TODO: اذكر المخاطر الرئيسية وطرق التعامل معها.
\begin{table}[H]
\centering
\caption{المخاطر الرئيسية}
\renewcommand{\arraystretch}{1.3}
\begin{tabularx}{\textwidth}{C{1cm} L{4cm} C{2cm} C{2cm} X}
\toprule
\textbf{رقم} & \textbf{المخاطرة} & \textbf{الاحتمال} & \textbf{التأثير} & \textbf{إجراء التخفيف} \\
\midrule
1 & تأخر المورد في التسليم & متوسط & عالي & تحديد مورد بديل مسبقاً \\
2 & تغير المتطلبات & عالي & متوسط & إدارة التغيير عبر عملية رسمية \\
3 & نقص الموارد البشرية & منخفض & عالي & تخطيط احتياطي للموارد \\
\bottomrule
\end{tabularx}
\end{table}

% ============================================================
% 9. COMMUNICATION PLAN
% ============================================================
\section{خطة التواصل}
\label{sec:communication}

% TODO: حدّد خطة التواصل مع أصحاب المصلحة.
\begin{table}[H]
\centering
\caption{خطة التواصل}
\renewcommand{\arraystretch}{1.3}
\begin{tabularx}{\textwidth}{L{3.5cm} L{3cm} C{2.5cm} X}
\toprule
\textbf{الجمهور} & \textbf{نوع التواصل} & \textbf{التكرار} & \textbf{المحتوى} \\
\midrule
الراعي الرئيسي & اجتماع & شهري & تقدم المشروع والميزانية والمخاطر \\
فريق المشروع & اجتماع & أسبوعي & المهام والتقدم والمعوقات \\
أصحاب المصلحة & تقرير & شهري & ملخص التقدم والمخرجات \\
المستخدمون & بريد إلكتروني & حسب الحاجة & التحديثات والإشعارات \\
\bottomrule
\end{tabularx}
\end{table}

\vspace{2em}

\noindent\rule{\textwidth}{0.4pt}

\vspace{0.5em}

\noindent
\begin{tabularx}{\textwidth}{X X}
\textbf{توقيع مدير المشروع:} & \textbf{توقيع الراعي الرئيسي:} \\[2em]
\rule{6cm}{0.3pt} & \rule{6cm}{0.3pt} \\
\end{tabularx}

\end{document}
